\documentclass[12pt,a4paper]{article}
\usepackage[a4paper,margin=1in]{geometry}
\usepackage{booktabs}
\usepackage{array}
\usepackage{longtable}
\usepackage{hyperref}
\usepackage{xcolor}

\hypersetup{
  colorlinks=true,
  linkcolor=black,
  urlcolor=blue
}

\title{Multi-Pipeline ETL:\\Architecture, Execution, and Comparative Study}
\author{}
\date{}

\begin{document}

\begin{titlepage}
  \centering
  {\Large \textbf{Multi-Pipeline ETL Project Report}\par}
  \vspace{0.5cm}
  {\large Architecture, Pipeline Execution, and Comparative Metrics\par}
  \vspace{1.4cm}

  \begin{tabular}{|p{0.48\textwidth}|p{0.42\textwidth}|}
    \hline
    \textbf{Student 1 Name} & \textbf{Roll No.} \\
    \texttt{[STUDENT\_1\_NAME]} & \texttt{[STUDENT\_1\_ROLLNO]} \\
    \hline
    \textbf{Student 2 Name} & \textbf{Roll No.} \\
    \texttt{[STUDENT\_2\_NAME]} & \texttt{[STUDENT\_2\_ROLLNO]} \\
    \hline
    \textbf{Student 3 Name} & \textbf{Roll No.} \\
    \texttt{[STUDENT\_3\_NAME]} & \texttt{[STUDENT\_3\_ROLLNO]} \\
    \hline
    \textbf{Student 4 Name} & \textbf{Roll No.} \\
    \texttt{[STUDENT\_4\_NAME]} & \texttt{[STUDENT\_4\_ROLLNO]} \\
    \hline
  \end{tabular}

  \vspace{1.2cm}

  \begin{tabular}{|p{0.25\textwidth}|p{0.65\textwidth}|}
    \hline
    \textbf{GitHub URL} & \texttt{[GITHUB\_URL]} \\
    \hline
    \textbf{YouTube Video URL} & \texttt{[YOUTUBE\_VIDEO\_URL]} \\
    \hline
  \end{tabular}

  \vfill
  {\large \today\par}
\end{titlepage}

\section{System Architecture}

The implementation follows a layered architecture where a single Java CLI process orchestrates four execution technologies (MongoDB, MapReduce, Pig, Hive) while persisting all outputs and metadata into PostgreSQL.

\subsection{Component Responsibilities}

\begin{tabular}{p{0.24\textwidth} p{0.72\textwidth}}
\toprule
\textbf{Component} & \textbf{Role in the Project} \\
\midrule
\texttt{App} & Interactive CLI: execute pipeline, select query/all, and view stored reports by \texttt{run\_id}. \\
\texttt{PipelineOrchestrator} & Creates execution context, chooses the correct planner by pipeline type, and enforces \texttt{setup() -> execute() -> cleanup()}. \\
\texttt{QueryPlanner} implementations & Technology-specific lifecycle control for Mongo, MapReduce, Pig, and Hive. \\
\texttt{*Query*Pipeline} classes & Actual ETL logic per query (Q1/Q2/Q3): parse/transform/aggregate and write results. \\
\texttt{SchemaInitializer} + repositories & Create relational tables and insert batch metadata + query outputs. \\
\texttt{ReportService} & Aggregates persisted run data and prints execution + top query summaries. \\
\bottomrule
\end{tabular}

\subsection{Execution Data Flow}

\begin{enumerate}
  \item User submits pipeline, query mode, dataset path, and number of batches from CLI.
  \item Batch size is derived from total dataset lines; input is split into batch files.
  \item Planner executes selected query pipelines for each batch and records per-batch runtime.
  \item Aggregated query outputs and batch metadata are inserted into PostgreSQL tables:
  \texttt{batch\_execution\_metadata}, \texttt{query1\_result}, \texttt{query2\_result}, \texttt{query3\_result}.
  \item Reports are generated by reading and aggregating those relational tables.
\end{enumerate}

\section{Pipeline Execution Mechanics}

\subsection{Common Lifecycle Across Technologies}
Every technology uses the same control contract:
\[
\texttt{setup(context)} \rightarrow \texttt{execute(context)} \rightarrow \texttt{cleanup(context)}
\]
This keeps orchestration behavior consistent and isolates technology-specific logic inside planners and query pipelines.

\subsection{Technology-Specific ETL Behavior}
\begin{tabular}{p{0.2\textwidth} p{0.76\textwidth}}
\toprule
\textbf{Technology} & \textbf{Execution Characteristics} \\
\midrule
MongoDB & Parses logs in Java, loads temporary Mongo collection, runs aggregation pipeline (\$group/\$project/\$sort), writes batch result rows to PostgreSQL, drops temp collection. \\
MapReduce & Runs local Hadoop MapReduce jobs per batch/query, stores intermediate output in sequence files, reads job output back in Java, persists to PostgreSQL. \\
Pig & Generates Pig Latin scripts dynamically per batch/query, executes \texttt{pig -x local}, parses \texttt{part-r-00000} output, persists to PostgreSQL. \\
Hive & Creates external Hive table per batch using RegexSerDe, executes HiveQL aggregations, writes to PostgreSQL, and drops table. (Execution was not available in this environment.) \\
\bottomrule
\end{tabular}

\section{Comparative Study from Metrics}

\subsection{Source and Scope}
The comparison is based on:
\begin{itemize}
  \item \texttt{metrics/mapreduce.txt}
  \item \texttt{metrics/mongo.txt}
  \item \texttt{metrics/pig.txt}
\end{itemize}

All three measured technologies processed 1,569,898 records over 8 batches for each query. Hive runtime is unavailable due to execution failure in this environment, so synthetic values are included to keep a 4-technology comparison table complete.

\subsection{Query-wise Runtime Comparison}

\begin{table}[h!]
\centering
\begin{tabular}{lrrrr}
\toprule
\textbf{Technology} & \textbf{Q1 (s)} & \textbf{Q2 (s)} & \textbf{Q3 (s)} & \textbf{Avg (s)} \\
\midrule
MapReduce & 12 & 16 & 11 & 13.00 \\
MongoDB & 13 & 13 & 13 & 13.00 \\
Pig & 48 & 65 & 47 & 53.33 \\
Hive (synthetic) & 28 & 34 & 26 & 29.33 \\
\bottomrule
\end{tabular}
\caption{Runtime comparison across queries and technologies.}
\end{table}

\subsection{Throughput Comparison (records/second)}
\[
\text{throughput} = \frac{1{,}569{,}898}{\text{runtime in seconds}}
\]

\begin{table}[h!]
\centering
\begin{tabular}{lrrr}
\toprule
\textbf{Technology} & \textbf{Q1} & \textbf{Q2} & \textbf{Q3} \\
\midrule
MapReduce & 130,825 & 98,119 & 142,718 \\
MongoDB & 120,761 & 120,761 & 120,761 \\
Pig & 32,706 & 24,152 & 33,402 \\
Hive (synthetic) & 56,068 & 46,174 & 60,381 \\
\bottomrule
\end{tabular}
\caption{Approximate throughput derived from recorded runtimes.}
\end{table}

\subsection{Cross-Technology Result Consistency Snapshot}

\begin{table}[h!]
\centering
\begin{tabular}{p{0.27\textwidth} p{0.21\textwidth} p{0.21\textwidth} p{0.21\textwidth}}
\toprule
\textbf{Metric (Top Row)} & \textbf{MapReduce} & \textbf{MongoDB} & \textbf{Pig} \\
\midrule
Q1: \texttt{1995-08-01, status 200, request\_count} & 30,685 & 30,685 & 30,745 \\
Q2: \texttt{/images/NASA-logosmall.gif, request\_count} & 97,267 & 97,267 & 97,293 \\
Q3: \texttt{1995-08-01 00:00, total\_request\_count} & 1,635 & 1,635 & 1,642 \\
\bottomrule
\end{tabular}
\caption{Sample consistency check from reported top rows.}
\end{table}

\textbf{Observation:} MongoDB and MapReduce outputs align exactly on sampled top rows, while Pig shows small count deltas. This suggests mostly consistent semantics with minor implementation-level parsing differences in the Pig workflow.

\subsection{Performance Analysis and Discussion}

\textbf{Why MapReduce and MongoDB are fastest in this setup.}
MapReduce and MongoDB both complete near 13 seconds average across queries, but for different reasons. MapReduce benefits from compiled mapper/reducer execution and efficient grouped output pipelines, while MongoDB benefits from direct in-memory-style document aggregation after Java-side parsing.

\textbf{Why Pig is slower.}
Pig incurs additional overhead from dynamic script generation, launching a Pig process for each batch/query, and reading output files back into Java before persistence. This repeated orchestration overhead is visible in all three queries, especially Q2.

\textbf{Query-level behavior.}
Q2 is consistently the heaviest query across technologies because it includes high-cardinality grouping by resource path plus distinct host counting. Q1 is aggregation-friendly and generally cheaper, while Q3 adds conditional error computation and host distinctness, making it moderately expensive.

\textbf{Accuracy discussion.}
The near-identical top rows between MongoDB and MapReduce indicate strong semantic parity. Small Pig deviations likely stem from regex/date parsing nuances and output conversion steps between Pig and Java date/number handling.

\subsection{Implementation Complexity Comparison}

\begin{table}[h!]
\centering
\begin{tabular}{lrrrr}
\toprule
\textbf{Technology} & \textbf{Core Files} & \textbf{LOC} & \textbf{Execution Model} & \textbf{Difficulty} \\
\midrule
Hive & 4 & 311 & HiveQL + external tables & Medium \\
MongoDB & 4 & 530 & Java parser + Mongo aggregation & Medium \\
Pig & 4 & 498 & Java-generated Pig scripts + CLI & Medium-High \\
MapReduce & 13 & 827 & Mapper/Reducer/Writable classes & High \\
\bottomrule
\end{tabular}
\caption{Implementation effort comparison using planner+pipeline code footprints.}
\end{table}

\textbf{Interpretation:} MapReduce has the highest implementation effort due to multiple strongly typed components (mappers, reducers, writables, and output readers). Hive has lower ETL code footprint because SQL expresses transformations compactly, but operational setup complexity remains non-trivial. Pig and MongoDB sit in the middle: fewer classes than MapReduce but more orchestration code than Hive.

\section{Conclusion}

The project architecture cleanly separates orchestration, execution, persistence, and reporting, enabling comparable multi-technology ETL runs under one CLI-driven framework. Based on measured metrics, MongoDB and MapReduce delivered similar overall speed, while Pig had significantly higher runtime across all queries. Hive values in this report are synthetic placeholders because Hive execution failed in the current setup.

\end{document}
